% Created 2026-04-09 qui 10:46
% Intended LaTeX compiler: xelatex
\documentclass[11pt]{article}
\usepackage{graphicx}
\usepackage{longtable}
\usepackage{wrapfig}
\usepackage{rotating}
\usepackage[normalem]{ulem}
\usepackage{capt-of}
\usepackage{hyperref}
\author{Nícolas Rosenthal Dal Corso}
\date{\today}
\title{Laboratório 2: Limpeza e Transformação de Dados}
\hypersetup{
 pdfauthor={Nícolas Rosenthal Dal Corso},
 pdftitle={Laboratório 2: Limpeza e Transformação de Dados},
 pdfkeywords={},
 pdfsubject={},
 pdfcreator={Emacs 30.2 (Org mode 9.7.11)}, 
 pdflang={English}}
\begin{document}

\maketitle
\tableofcontents

\newpage
\section{Introdução}
\label{sec:org3dd683a}
O objetivo deste trabalho foi transformar um dataset com problemas de qualidade em um conjunto de dados limpo, consistente e adequado para análise.

O dataset original apresentava problemas típicos de dados reais, incluindo:
\begin{itemize}
\item valores duplicados
\item inconsistência em texto (nomes e categorias)
\item idades inválidas
\item valores ausentes
\item possíveis outliers em renda
\end{itemize}
\section{Inspeção Inicial}
\label{sec:orgb004d26}
A análise inicial mostrou:
\begin{itemize}
\item duplicatas exatas (ex: registros repetidos de clientes)
\item nomes com capitalização inconsistente (ex: "carla")
\item idades inválidas (ex: 200 anos)
\item possíveis inconsistências em categorias como cidade e gênero
\end{itemize}
\section{Etapas de Limpeza}
\label{sec:orgd5a4fee}

\subsection{Remoção de Duplicatas}
\label{sec:orgb5b4b41}
Registros duplicados foram removidos para evitar viés estatístico.
\subsection{Padronização de Texto}
\label{sec:orgf04be0b}
As colunas textuais foram normalizadas:
\begin{itemize}
\item nomes capitalizados corretamente
\item categorias padronizadas (ex: gênero e cidade)
\end{itemize}
\subsection{Tratamento de Idades Inválidas}
\label{sec:org58a527d}
Idades irreais (ex: > 100 ou negativas) foram substituídas por valores ausentes (NaN), evitando distorções na análise.
\subsection{Conversão e Padronização de Datas}
\label{sec:org24658e5}
A coluna \texttt{signup\_date} foi convertida para formato datetime, tratando múltiplos formatos de entrada.
\subsection{Tratamento de Valores Ausentes}
\label{sec:orgd3fdc56}
Foi adotada a seguinte estratégia:
\begin{itemize}
\item variáveis numéricas: preenchidas com a mediana
\item variáveis categóricas: preenchidas com "Desconhecido"
\end{itemize}
\section{Transformações Realizadas}
\label{sec:org6e8c689}

Foram criadas novas variáveis para enriquecer o dataset:

\begin{itemize}
\item \texttt{signup\_month}: mês de cadastro extraído da data
\item \texttt{income\_minmax}: normalização da renda (escala 0-1)
\item \texttt{income\_log}: transformação logarítmica para reduzir assimetria
\item \texttt{segment\_encoded}: codificação numérica da variável categórica
\item \texttt{purchases\_per\_age}: razão entre compras e idade
\item \texttt{income\_per\_purchase}: renda média por compra
\end{itemize}

Além disso, a variável categórica \texttt{city} foi transformada usando one-hot encoding, gerando colunas como:
\begin{itemize}
\item city\_Porto Alegre
\item city\_São Paulo
\item city\_Rio de Janeiro
\item city\_Curitiba
\item city\_Desconhecido
\end{itemize}
\section{Outliers}
\label{sec:orgca111cc}
Foi identificado um valor de renda potencialmente extremo.

Optou-se por manter o valor, aplicando transformação logarítmica (\texttt{income\_log}), pois:
\begin{itemize}
\item a remoção poderia descartar informação relevante
\item a transformação reduz o impacto do outlier sem perder o dado
\end{itemize}
\section{Dataset Final}
\label{sec:org5ccc57a}
O dataset final contém:
\begin{itemize}
\item dados limpos e consistentes
\item variáveis transformadas relevantes
\item codificação adequada para modelos de machine learning
\end{itemize}

O conjunto final possui 21 colunas, incluindo variáveis derivadas e codificadas.
\end{document}
